% TODO proof-reading
\subsection{ビットフィールド}

すべてのビット演算は他の\ac{ISA}と同様に動作します：

\begin{lstlisting}[style=customjava]
	public static int set (int a, int b) 
	{
		return a | 1<<b;
	}

	public static int clear (int a, int b) 
	{
		return a & (~(1<<b));
	}
\end{lstlisting}

\begin{lstlisting}
  public static int set(int, int);
    flags: ACC_PUBLIC, ACC_STATIC
    Code:
      stack=3, locals=2, args_size=2
         0: iload_0       
         1: iconst_1      
         2: iload_1       
         3: ishl          
         4: ior           
         5: ireturn       

  public static int clear(int, int);
    flags: ACC_PUBLIC, ACC_STATIC
    Code:
      stack=3, locals=2, args_size=2
         0: iload_0       
         1: iconst_1      
         2: iload_1       
         3: ishl          
         4: iconst_m1     
         5: ixor          
         6: iand          
         7: ireturn       
\end{lstlisting}

\TT{iconst\_m1}は$-1$をスタックにロードします。これは\TT{0xFFFFFFFF}数と同じです。

\TT{0xFFFFFFFF}とのXORは、すべてのビットを反転する効果と同じです
（\myref{XOR_property}）。

すべてのデータ型を64ビット\IT{long}に拡張してみましょう：

\begin{lstlisting}[style=customjava]
	public static long lset (long a, int b) 
	{
		return a | 1<<b;
	}

	public static long lclear (long a, int b) 
	{
		return a & (~(1<<b));
	}
\end{lstlisting}

\begin{lstlisting}
  public static long lset(long, int);
    flags: ACC_PUBLIC, ACC_STATIC
    Code:
      stack=4, locals=3, args_size=2
         0: lload_0       
         1: iconst_1      
         2: iload_2       
         3: ishl          
         4: i2l           
         5: lor           
         6: lreturn       

  public static long lclear(long, int);
    flags: ACC_PUBLIC, ACC_STATIC
    Code:
      stack=4, locals=3, args_size=2
         0: lload_0       
         1: iconst_1      
         2: iload_2       
         3: ishl          
         4: iconst_m1     
         5: ixor          
         6: i2l           
         7: land          
         8: lreturn       
\end{lstlisting}

コードは同じですが、64ビット値で動作する\IT{l}プレフィックスを持つ命令が使用されます。

また、関数の2番目の引数は依然として\IT{int}型で、その中の32ビット値を64ビット値にプロモートする必要がある場合、\TT{i2l}命令が使用されます。これは本質的に\IT{integer}型の値を\IT{long}型に拡張します。